\section{Two-port bound in oblique sections}\label{L:part9}

We give the power upper bound from \eqref{L:eq7}. We need it with both types of triangular rhombi (\(q,q^2\) as cell ratios), i.e., up to common rotation the spatial rapidities are \(1,q\). For positive integers \(m,K_0\), take \(m\) unmarked sites between the two marks on the first side of the slice, and \(K_0m\) on the other. On both parts take the same proportions \(1-\theta,\theta\) of the rapidities \(1,q\), \(0\le\theta\le1/2\), with actual integer counts. Mark rapidities \(a,b\) can each have either value. The bound we will show (large \(m\ge K_0^{1/4}\)) is
\begin{equation}
 0\le C_{ab}\le C m^{-4/3}K_0^{1/3}.\label{L:eq8}
\end{equation}

These are physical evaluations of the vacuum formulas. Take a row boundary of forward lattice side steps in two directions \(e_0,e_1\) of consecutive angles \(0,\pi/3\) in the triangular lattice; successive rows translate by \(e_1-e_0\). This places exactly the triangular splittings for \(q,q^2\), respectively, in the rhombi. The horizontal period is the sum of the steps. Half-cylinder vacuum convergence still holds here. In fact the lifted region has the forward zigzag as boundary. Exhaust the region by a finite zigzag segment and its distant parallel translate, with straight connecting sides parallel to the row translation. For a midport source on the base the lift of counterclockwise tangent angle, starting there with angle \(\alpha_s\), ranges in \([0,7\pi/3]\) until returning to the source. The turn of any simple arc to another boundary port of lifted angle \(\alpha\) is \(\alpha-\alpha_s-\pi\) (close it along the boundary). Thus the complex boundary flux argument with turn multiplier \(3/8\), rotated by a fixed phase at the source, has strictly positive real coefficients bounded below. Indeed the outer-source identity sums \(X^{|\omega|}e^{i(3/8){\rm turn}}\) to 1 by the two extensions at each step and reversal of collision detours from outside. It gives a uniform bound on the mass to base ports. The same exhaustion bound works on the other side by symmetry. This suffices for the previous convergence and accessibility proof (all single routes positive, double weight the product of allowed separate triangle routes when nonzero). The rational normalized vacuum coincides with the convergent actual one by transfer invariance and uniqueness wherever its denominator does not vanish; nonvanishing of \(Q\) also follows from the next comparison.

For \(Z\) such a list define
\[
\begin{gathered}
 c=3/4,\quad b_d=\left|\frac{h(1,q^d)}{1-q^d}\tan\frac{\pi d}{8c}\right|,
 \quad b_0=\frac{2+\sqrt2}{2c},\\
 {\cal D}(Z)=(2c)^{-|Z|/2}\prod_{i<j} b_{k_j-k_i},\qquad z_j=q^{k_j},\ k_j\in\{0,1\}.
\end{gathered}
\]
We have \(|Q(Z)|\ge c_1{\cal D}(Z)\). Use the positive Fourier weight comparison \(\mu\ge\mu_0\) of the calibration, now before confluence with distinct ordered real angles in a range slightly larger than \([0,\pi/4]\). By expanding exterior powers, the Fourier skew Pfaffian is (up to common phase and factorial) an integral over positive frequencies \(\nu_j\) with columns \(e^{-\nu_j\arg z},e^{\nu_j\arg z}\) in a determinant, and column 1 in augmented size, against \(\prod(\mu(\nu_j)/\nu_j)d\nu_j\). Differences at zero are taken as usual (absolutely integrable). The determinants all have fixed sign after ordering frequencies: real exponentials with distinct real exponents on an ordered real set form a nonsingular matrix, by successive differentiation/Rolle after dividing by one exponential. Hence replacement by \(\mu_0\) only decreases the modulus. The resulting skew kernel is \((2c)^{-1}\tanh(\log(x/y)/(2c))\), by the transform already used. Its Pfaffian gives a factor \((2c)^{-\lfloor |Z|/2\rfloor}\) times the product of individual hyperbolic tangents, by the identity for \((u_i-u_j)/(u_i+u_j)\) (compare residues by pair deletion, then the remaining constant by antisymmetry; use last \(u_j=0\) for augmentation). Multiply by \(F\) and take confluence.

We will use the same expansion for the two finite observables
\[
 F_Z^{(2)}=Q(Z)^2C_{ab},\qquad F_Z^{(0)}=Q(Z)^2.
\]
The unmarked expansion, together with the positive comparison just proved,
will determine the size of $|Q(Z)|$. The marked expansion will then be
calibrated by the elementary physical bound $C_{ab}\le C_a$, where $C_a$
is the mass through one port. This separates the normalization from the
decay between the marks.

Here are the size conventions used in the two calculations.
\begin{center}
\begin{tabular}{c|c|c|c}
 observable & regular family counts & $s_1$ & $K_0$\\\hline
 $F_Z^{(2)}$ & $m$ and $K_0m$ & $m$ & positive integer\\
 $F_Z^{(0)}$ & each angle class split nearly equally & $|Z|$ & $1$
\end{tabular}
\end{center}
In the marked case $|Z|=(K_0+1)m+2$, and the later-angle fraction is
$\theta$ in both regular families. Each mark can have rapidity $1$ or $q$.
For the unmarked case use the fugacity-zero identity~\eqref{L:eq2};
partition each angle class as equally as possible between the families,
with order at our disposal. Orient the two angles so that the total later
fraction is at most $1/2$, using inversion and common scaling if necessary.
These operations preserve $|Q|$ and ${\cal D}$.

Write $p_*=0,2$ for the number of marks. In the finite color and mark
weights use the conjugate convention $q\mapsto q^{-1}$ and
$\phi_+=\phi_-=-\pi/4$; the physical sites still have later rapidity $q$.
Indeed the glued mass is real on nearby unit-circle arrays and invariant,
as a rational function, under simultaneous rapidity inversion by
upper/lower reflection. Its coefficients can therefore be conjugated.

The composition parameters are the two angle fractions in each regular
family in the marked case, and all separate regular multiplicities divided
by $s_1$ in the unmarked case. We allow them to vary in fixed small complex
neighborhoods of the indicated limiting compositions, including a zero
later fraction. The final physical bound uses identical real later
fractions in the two marked families. Angular boundary limits are taken
by blocks, one for each angle and field kind, in a fixed order within an
angle.

There is also a finite phase index. For each fixed root list, its
orders and lattice indices, the part of the residual phase depending on
regular multiplicities is a rational quadratic polynomial in those
multiplicities. Here the physical source angles are $0,1$; the generic
root and counter-field angles and the bounded imaginary period shift
remain contour data. Express the axis index as $f=1/4+n$ with integer
$n$ and write any binomial terms in ordinary monomials. Before summing root occurrences, choose one integer $D_{\rm ph}$
clearing the finitely many underlying rational charge and bilinear-form
coefficients in these monomial expressions. Each root-list phase
polynomial is an integer combination of those coefficients. Thus the
same $D_{\rm ph}$ works for every list and order, independently of
$K_0$ and $J$. Fix the regular multiplicities
modulo $2D_{\rm ph}$. Changing them by this modulus changes each such
polynomial by an even integer. Their phase dependence is therefore fixed
on these classes. This does not freeze the phases depending on root lists,
orders or continuous contour data.

\begin{proposition}[Two-scale expansion]\label{L:prop:two-scale}
For an admissible integer array of either of the two kinds specified
above, put
\[
 G_Z^{(p_*)}=\frac{F_Z^{(p_*)}}{{\cal D}(Z)^2(K_0s_1)^{3/8}}.
\]
For each $J>0$ and each of the fixed composition neighborhoods, there
is a threshold $s_J$ independent of $K_0\ge1$ such that, on every fixed
phase class and every admissible integer array with $s_1\ge s_J$,
this quantity has a finite approximation
\begin{equation}
 G_Z^{(p_*)}=
 \sum_\gamma s_1^{-\gamma}P_\gamma(\log s_1;K_0,\text{parameters})
       +O_J(K_0^{C_2}s_1^{-J}).\label{L:eq9}
\end{equation}
The formula is asserted at admissible integer arrays. Its coefficient
functions extend holomorphically to the composition neighborhoods above and
are bounded by $C_JK_0^{C_2}$. One exponent $C_2$ works for every accuracy $J$;
the finite rate and degree lists may depend on $J$ but not on $K_0$.
In the marked case $\gamma\ge3/2$, with no logarithm at equality.
In the unmarked case $\gamma\ge1/6$, with no logarithm at equality,
and all other rates are at least $1/2$ higher.
\end{proposition}

We first construct the trace representing these finite quantities and
remove its extensive scalar factors. The proof of the proposition then
separates an edge of length $6\log K_0$ from the free propagation length
$6\log s_1$. Only the latter enters the accuracy-dependent convolutions.
This is what keeps the exponent $C_2$ fixed.

\subsection{Trace and bulk normalization}\label{L:step14}

Use all the doubled data of the enclosing-loop calculation. Introduce one charge \(H_\diamond\) in the common sector of that lattice by
\[
 H_\diamond D_j^\sigma=\mathbf1_{j=6}-\mathbf1_{j=3}.
\]
For each mark use a field of charge \(Y=H_\diamond\) at the mark, and place one counter-field (for both) of charge \(-2H_\diamond\) near the outer seam, say at \(u=L/2-c_s+2\pi i d_c\) for large period, \(c_s>0\) fixed small and \(0<d_c<1\) generic. Only use these extra fields in the marked case. On probe fields use translation exponent \(\mathcal B(Y,P)-2SY\cdot P\); on other fields use \(\mathcal B\) as before. The probe modification has zero total coefficient on the initial momentum (leaving only inter-field translation phases). The vector \(H_\diamond\) has \(h_o H_\diamond=v H_\diamond=0,\ hH_\diamond=3/4\). Its modified exchange sign with a root is ordinary (by \(S+S^t=1\)); and products with \(D_j^\sigma\) of indices \(7,0,1,2\) vanish. Thus probes are unaffected by the color conjugators and commute with all root operations used here.

Temporarily separate all sites \(z=e^{u/8}\) with \(u=x+2\pi i d\) near the indicated angles \(d=0,1\), in increasing angular order within a range strictly less than 1; take regular fields in increasing real order also, and all probes later than them in real order. Insert \(a^\sigma\) at regular sites, conjugated by the enclosing-loop \(g\)'s. Divide the resulting trace by the following uncolored prescription: take the \(O\) scalar channel on uncolored \(a^\sigma\) and all \(Y\), with its prefactor, self constants including field powers, and direct short factors only, omitting also the direct factors to the counter-field. Use the phase (units of \(\pi\))
\[
 -\sum_i\Delta_{\beta_i}-\sum_{i<j}[\mathcal B(\beta_i,\beta_j)+2SY_j\cdot\beta_i]
 -\sum_{i<j}{\rm sgn}(x_j-x_i)\beta_i T_O\beta_j
\]
in angular order, \(Y_j=0\) at regular sites; use for the kinetic factor
\(\prod_{x_i<x_j}\exp((h_o\beta_i)(h_o\beta_j)(u_j-u_i)/(2H_o))\).
This is a formal normalization, not requiring \(h_o\)-neutrality. Finally multiply by
\[
 \prod_{{\rm reg}\ i<j}D_{\sigma_i\sigma_j}
 \prod_{{\rm reg}\ i}(\sqrt\rho\,e^{i\phi_{\sigma_i}})
\]
and, if marked, also by \(\rho ab\prod_{z\,{\rm reg}} f_+'(a,z) f_+'(b,z)\), with \(f'\) the mark polynomial in conjugate convention.

The period constant coefficient (\(p=e^{-L/8}\), allowing half powers) of this prescription is \(F_Z^{(p_*)}\). Indeed expansion in original colors gives exactly the phase step computation at \(O\) from \eqref{L:eq2}: only neutral colors occur and at \(f=0\) the kinetic term is the pair expression above on them by neutrality. Thus regular-regular ratios, self factors and color coefficients give the same pair and site factors as before. For a color \(b'\) and outer probe \(Y\), the direct relative factor against the uncolored site is
\[
  \prod_j(1-q^{-j}z/z_{\rm probe})^{YO^j(b'-a^\sigma)}.
\]
Mixed phase from the indicated pair equals \(-2b'Y\) if the color is earlier in angle, and zero otherwise, hence no relative phase. Our row on the two color changes \(-D_0,-D_0-D_1\) thus supplies \(f_c'/f_+'\) at each mark, and all counter-field direct factors tend to 1. Image factors are regular additional series and nonzero \(f\)'s have \(p^{f^2}\).

Removing the conjugators gives the \(M_*\) channel with precisely the same root species \eqref{L:eq3} and coefficients as before. All root-circle fields there commute through uncolored regular fields (orthogonal products), so can be interspersed with them keeping the indicated order among the circles. The total charge in either channel is unchanged by probes, and the only new translation effects besides the ordinary \(\mathcal B\) formula are \(-2\sum_{i<j}SY_j\cdot\beta_i\). Thus all the previous cut/phase tests work (integral root products on the dual lattice); at \(M_*\) require \(\sum h\beta=0\), with the same possible parity obstruction and step \(h\). Use \(f\in f_0+\mathbb Z,\ f_0=1/4\). All residual phases from the cut, translations and order signs have rational data, depending as regards regular multiplicities only on congruence classes. Formal properness, oscillator contraction and continuation at finite real period are justified exactly by the doubled estimates: external charges have fixed bounded multiplicities during those identities, trace reordering uses probes commuting as above, cut support and quadratic time-average tests change only by bounded charge translates, and root counts themselves have neutral total \(h\)-charge. The bounded axis shifts from probes can be included using a portion of kinetic margin in the bin estimate, at cost exponential in total real length. Separated root and external angles then give absolute summability. For reaching a limiting counter position one can first continue with it strictly inside the interval then move it at large period at its nonsingular angle. Near the seam bend use a constant bend at that position.

Here are contraction data. For regular \(a^\sigma\) on \(M_*\) we have \(a^\sigma(M_*^2-M_*+I)=0\) as contractions. Its row \(a^\sigma M_*^j a^\tau\), or \(a^\sigma M_*^j H_\diamond\), is \((4,2,-2,-4,-2,2)/3\). On the five root columns in order \eqref{L:eq3}, \(a^\sigma M_*^j r_i^\tau\) vanishes for opposite type and for the same type is \((1,0,-1,0,-1)_i\) times
\[
 c_j'=(0,-1,-1,0,1,1)_j .
\]
On \(O\) the rows against \(a^\sigma,a^{-\sigma},H_\diamond\) respectively are
\[
 (4,2,1,-1,-2,-1,1,2)/3,\quad
 (4,2,-2,-4,-2,-4,-2,2)/3,\quad
 (4,2,1,-1,-2,-4,-2,2)/3 .
\]
Also \(H_\diamond M_*^{-1}\delta^\sigma=1\). For example to check by multiplication, on sum/difference use cyclic coefficient columns with shift \(V\), metric circulant \(C\) the sum/difference of the two rows given before, regular column \((I-V+V^2)^{-1}e_0\). Multiply by \(V(I-e_7 e_7^{\mathsf T}C)\) for \(M_*\); the common coefficient column of \(H_\diamond\) (same in both families) is \((7,-7,-5,-19,-17,17,19,5)^{\mathsf T}/48\).

At confluence (including angular boundary contact), divide each direct central pair factor in the \(M_*\) expression by its uncolored \(O\) counterpart. Indeed \(K_U(Y'+2\pi i d)=U^dK_U(Y')\), \(d=0,\pm1\), and the singular exponent \(\gamma_i U^d\gamma_j\) agrees on both sides for uncolored central fields (all orthogonal to both \(\delta\)'s). Thus there is a finite constant ratio in increasing real order. When rows \(\gamma_i U^j\gamma_j=A(j)\) have finite period \(l\), the modulus of the constant of that individual short factor after its power is removed is
\[
 T(A,d)=l^{-A(d)}\prod_{j=1}^{l-1}|2\sin(\pi j/l)|^{A(j+d)},
\]
and its phase rational here by the directed log formula. With either two regular sites or regular then mark, multiplication of \(T(M_*,d)/T(O,d)\), \(d=k_j-k_i\), by respectively \(|D_{\sigma_i\sigma_j}(z_i,z_j)|\) or \(|f_+'(z_j,z_i)|\) gives precisely \(b_d^2\). These identities follow by inserting the displayed 6- and 8-rows for \(d=-1,0,1\), with \(|D_+|=|(1+q^{2d})(1+q^{2d}+\sqrt2 q^d)|,\ |D_-|=1,\ |f_+'|=|(q^d-q^2)(q^d-q^3)|\). Self ratios similarly give \(\sqrt\rho\, d_{M_*}(a^\sigma)/d_O(a^\sigma)=1/(2c)\). Mark-only constants (also the counter self ratio) are bounded independently of size. Kinetic phases from the formal normalization at confluence are unit with rational data. Hence the whole extensive constant has modulus \({\cal D}(Z)^2\) up to fixed factors, and after dividing this out the dependence on regular multiplicities there is at most by congruence class. Field period powers cancel (root powers against measures). The determinant/prefactor ratio gives \(e^{L/32}\) times a convergent series at large \(\Re L\) as calculated earlier.

Initially choose first root layer at small \(d_s>0\), others at \(d_s+1/2,d_s+d_f<1\) (notation \eqref{L:eq3}). They stay clear of confluences. Absolute bin control permits taking confluence in the transformed ratio. On the old color side take just the constant coefficient in separation along generic fixed slopes: direct color ratios are meromorphic with bounded pole orders, other image series regular there and at \(p^{1/2}=0\), including the omitted counter factors as regular series. Thus this limit gives the claimed coefficient rule at confluence as well.


\subsection{Circle expansion at two scales}\label{L:step15}

\begin{proof}[Proof of Proposition~\ref{L:prop:two-scale}]
\textbf{Suppressive contours.}
Here is a proof with the two scales. Cauchy average with bounded \(\Im L\) (an interval of length \(32\pi\) for the expansion in \(p^{1/2}\)) and real half-period
\[
 \ell=\Delta_1+\lambda_1,\qquad \Delta_1=6\log K_0,\qquad \lambda_1=6\log s_1.
\]
Measure each real half arc by distance \(r'\) from the seam to the central fields at 0. Use the common small-slope bend supported in fixed seam widths as before. The first layer \(d=d_s\) should have a strictly suppressive unit profile \(c_j'\); the middle and last at \(d_s+1/2,d_s+d_f\) should have strictly suppressive negative unit profile (for each family's mixture). To see the choice, the unit log for rapidity 1 is
\[
 S_1(d;t)=\sum_j c_j'\log|1-t e^{2\pi i(d-j)/6}|,\qquad 0<t\le1 .
\]
It is odd and strictly increasing on \((-1,1)\): twice the log expression means the log of the ratio \((1-a'^2/4+a'^2 b'^2+\sqrt3 a'b')/(1-a'^2/4+a'^2 b'^2-\sqrt3 a'b')\), \(a'=2t/(1+t^2), b'=\sin(\pi d/3)\). The rapidity \(q\) uses \(d-1\). For a fixed common limiting later fraction \(0<\theta\le1/2\), take small \(0<d_s<1-d_f\). Then \((1-\theta)S_1(d;t)+\theta S_1(d-1;t)\) has the desired strict signs (negative at the first, positive at the other two), uniformly after division by \(t\), so also for both actual mixtures nearby. For limiting fraction zero take a small \textbf{negative} \(d_s\); now the main first profile is strictly negative, the other source log at \(d_s-1\) finite.

To justify the latter contour choice, work before bends at positive real period and first with sufficiently close but separated sources as above. Keep the phase prescriptions of the original angular order. For each fixed root list slide all its angles together by the required small amount, choosing the generic counter angle away from the swept layers modulo 1. In the angular contraction for \(i<j\) the position ratios enter oscillator factors
\[
 (1-\mathfrak q^n s_i/s_j)^{\beta_i M_*^{-n}\beta_j}\ (n\ge0),
 \qquad
 (1-\mathfrak q^n s_j/s_i)^{\beta_i M_*^n\beta_j}\ (n\ge1).
\]
Those involving a root have integral powers (\(M_*\) preserves the root lattice and its dual), as do its angular momentum powers. Root-root contacts within layers are nonnegative powers and their ratios are unchanged. The only possible new collisions involve \(\delta^\sigma\), with exponent zero against first-angle charges at direct contact, and exponent \(a^\tau M_*^{-1}\delta^\sigma=\mathbf1_{\sigma=\tau}\) or \(H_\diamond M_*^{-1}\delta^\sigma=1\) against later ones at translated contact. Thus there is no pole or branch obstruction: the integrand for a given angular momentum is analytic on scaling all root variables by a common factor across these radii (the remaining infinite products normally convergent), and its torus average is rotation invariant in the factor, hence constant. The momentum sum for the fixed list is Gaussian-convergent. Off real-part collisions, its radial scalar prescription continues to agree along this angle change by holomorphy, using the same real orders and phases (\(\Re Y>0\) throughout on all the corresponding short paths). Counts need only be re-summed at the two endpoint prescriptions, by the bin estimate, before taking confluence; each has all root contours clear of external contact there. The earlier small-bend estimates then give the continuation to complex periods used in the average.

A coupled root (nonzero profile against regulars) has source-to-root short multiplier in modulus bounded by
\[
 C\exp(-c_2 e^{v_i/6}),\qquad
 v_i=\begin{cases} r'_i-\Delta_1,&\text{size }s_1\text{ family},\\
                    r'_i,&\text{size }K_0s_1\text{ family}.
       \end{cases}
\]
For \(K_0=1\) use \(v_i=r'_i\) for both, with totals comparable to \(s_1\). This follows by the strict signs at \(t=e^{-(\ell-r')/6}\) and the displayed contraction rows, also for a root on the reverse side. More distant images cost a bounded factor, and within the bend this bound only needs bounded log estimates. Uncoupled roots (\(-D_2,\ \delta-D_1-D_2\)) change axis momentum by \emph{positive} integers everywhere.

Call \(V_1=\max(0,\max_{i\ {\rm coupled}}v_i)\). Exclude first \(V_1>\lambda_1/2\) with arbitrarily high exponential accuracy in \(\lambda_1\), at cost \(e^{C\Delta_1}\). Indeed the bin estimate for the remaining interactions, including probe costs, and retaining some kinetic fraction, bounds total volumes (without this penalty) by \(e^{C(\Delta_1+\lambda_1)}\): there are exponential-series factorials for the root species, group and measure/self costs a fixed-base exponential per particle, and the \(f\) sum costs \(O(1+\#\text{particles})\) by Gaussian tails. Probe short factors per moving particle are bounded since angles were taken strictly off contacts. Comparing kinetic stability of roots alone to actual kinetic energy uses a bounded shift on any interval and loses only \(C\) per unit length, by using a slightly larger fraction of actual energy. Regular self/central image terms cost \(O(1)\) in the log since only frequencies in \(\mathbb Z\pm1/6\) occur there against regular charges.

\par\medskip\noindent\textbf{Source expansion.}
Taylor expand all the remaining factors involving regular source multiplicities (bulk and phases already treated), and the determinant correction after removing \(e^{L/32}\), at fixed seam coordinates in \(s_1^{-1/2}\). The log terms of order zero in the source interactions are kept exponentiated. To detail the bounds, a near directed path of positive frequency \(m'\in\mathbb Z\pm1/6\), with \(n'\) extra periods, contributes for a coupled root with log magnitude controlled coefficientwise by a constant times
\[
 s_1^{-(6m'(1+2n')-1)}
 \, k' \exp[-m'(1+2n')\Delta_1+m'r'],\qquad k'=1\text{ or }K_0
\]
using that family's size; a far path has \(-m'r'\) instead. The only order-zero terms are first harmonics. They themselves obey the displayed localization, by the limit or its strict first-harmonic sign, also when the composition parameters are taken in a fixed small complex neighborhood. (Phases involving extensive integer data not in these harmonic logs are held fixed according to the size class, and \(h a^\sigma=0\).) The prefactor \(k'e^{-m'\Delta_1+m'r'}\) is bounded by \(e^{m'\max(0,v_i)}\). Thus the expansions converge with modulus still at most \(C\exp(-c_3 e^{v_i/6})\) per coupled root for dummy \(|s_1^{-1}|\le c_4 e^{-V_1/4}\) with sufficiently small \(c_4\); higher harmonics are geometric and cost an arbitrarily small multiple of \(1+e^{v_i/6}\) in log there. This concerns the interactions at fixed seam coordinates, before length integrations. Other terms with regulars are central images at integral periods (at most order \((K_0s_1)^2\) charge factor), or counter paths at odd half periods plus constant offsets (linear); they have bounded expansions on that disk. Determinant corrections use positive period powers corresponding to \(F_j\) above, i.e. powers generated by \(12/j\), \(j=1,2,3,6,8\), in \(s_1^{-1}\).

To keep track of the order dependence, let $P$ be the number of roots and let $A_P(z)$ be this complete
remaining source factor for one root configuration, including the scalar
source-image and determinant terms, with $z=s_1^{-1/2}$. Put
\[
 S_1'=\sum_{i\ {\rm coupled}}e^{v_i/6},\qquad
 r(V_1)=\sqrt{c_4}\,e^{-V_1/8}.
\]
The preceding disk estimate is
\[
 \sup_{|z|\le r(V_1)}|A_P(z)|\le C D_*^P e^{-cS_1'}
\]
with fixed $D_*,c>0$. Apply Cauchy's estimate once to the whole product:
\[
 |[z^n]A_P|\le Cc_4^{-n/2}e^{nV_1/8}D_*^P e^{-cS_1'}
              \le C_nD_*^P e^{-cS_1'/2}.
\]
For $V_1>0$ the last inequality uses $S_1'\ge e^{V_1/6}$; for $V_1=0$
it is immediate. Thus the Taylor order changes $C_n$, not the base
$D_*$ raised to the number of roots.

On $V_1\le\lambda_1/2$, the actual $|z|/r(V_1)$ is at most
$c_4^{-1/2}s_1^{-1/8}$, hence below $1/2$ for large $s_1$.
The degree-$N$ remainder, by the same Cauchy estimate and absorption,
is bounded by $C_Ns_1^{-(N+1)/2}D_*^Pe^{-cS_1'/2}$.
Integrating the previously established volume bound
$e^{C(\Delta_1+\lambda_1)}$ gives
\[
 C_NK_0^{6C}s_1^{6C-(N+1)/2}.
\]
Choose $(N+1)/2>J+6C$. This proves the source-truncation estimate
with a fixed exponent on $K_0$.

\paragraph{Which length is expanded.}
We now integrate a fixed coefficient of the source Taylor expansion.
The estimate above retains localization and a fixed exponential base
per root. Write $P$ for the root count in the remainder of the proof.

On each half-arc put an artificial node at $\Delta_1+R_0$, beyond the
bend and counter-field, and let $\Delta_1+R_j$ be the last of this node
and all coupled roots on that arc, $j=L,R$. This is the edge of that
arc. Localization gives exponential moments of $R_j$ of every fixed
order. The edge may be long because of $\Delta_1$, but its remaining
length up to the central fields is
\[
             \lambda_1-R_j.
\]
The source Taylor coefficients depend only on the edge data and the
fixed parameters. In particular they do not depend on how this
remaining length is divided into free gaps. Figure~\ref{L:fig:two-scale} records this separation.
\begin{figure}[htbp]
\centering\begin{tikzpicture}[x=1cm,y=1cm,>=Stealth]
\draw[very thick,blue!60!black] (0,0)--(4.4,0);
\draw[very thick,red!65!black] (4.4,0)--(9,0);
\foreach \x in {0,3,4.4,9} \draw (\x,-.1)--(\x,.1);
\node[below] at (0,-.12) {seam};
\node[below] at (3,-.12) {$\Delta_1$};
\node[below] at (4.4,-.12) {$\Delta_1+R_j$};
\node[below,align=center] at (9,-.12) {central fields\\$\ell=\Delta_1+\lambda_1$};
\draw[<->] (0,.45)--node[above] {edge}(4.4,.45);
\draw[<->] (4.4,.45)--node[above] {free length $\lambda_1-R_j$}(9,.45);
\end{tikzpicture}

\caption{One real half-arc. The edge ends at $\Delta_1+R_j$;
its complement has length $\lambda_1-R_j$.}
\label{L:fig:two-scale}
\end{figure}

Sort the roots after the edge and denote the free gaps by $d_i'\ge0$.
On each arc their sum is $\lambda_1-R_j$. Their kinetic factors are
$e^{-E(y_i)d_i'}$, where $E(y)=y^2/(2H)$ and $H=9/4$. The momenta
change by positive integers in increasing real order; reversing the
parameterization on the other arc reverses all these changes. Thus,
for every fixed $B$, the number of gaps with $E(y_i)\le B$ is bounded
independently of $P$. Imaginary kinetic phases use only the fixed
angles and the edge data.

\paragraph{Exposing the long gaps.}
Fix a sufficiently large $s_0$, and let $B$ be an energy threshold to
be chosen below. A gap is active when
$E(y_i)\le B$ and $d_i'\ge s_0$. Expand every short interaction whose
directed path crosses an active gap, including image paths. First use
the frequency series for its logarithm and then exponentiate. Each
resulting path label contributes a nonnegative shift $\sigma_i$ to
 every crossed active rate and may contribute a polynomial in its
length from Jordan differentiation. All positive frequencies belong
to $\tfrac16\mathbb Z$, so every polynomial contribution is accompanied
by a positive shift. Interactions on paths crossing no active gap are
independent of active lengths and remain unexpanded.

The absolute sum of these expanded logarithms, even after reserving
half their frequency decay, is at most
\[
       \varepsilon\Bigl(1+\sum_I K_I^2\Bigr),
\]
where $\varepsilon$ is arbitrarily small for large $s_0$ and $K_I$
counts roots in unit bins. Indeed every such path has length at least
$s_0$; exponential frequency decay, including the polynomial Jordan
factors, is summable by bins. The same estimate bounds the cost of
deleting these paths from the intact product. The doubled stability
bound therefore still applies. It applies as well if active lengths
are varied independently above $s_0$: use the resulting circle, which
may have unequal half-arcs, with the same edge bend and angular
offsets. This is the estimate that allows some active lengths to be
set equal to $s_0$ below.

\paragraph{The bounds needed before extending coefficients.}
Fix the bend bound, distant-path margin, per-root exponential base and
retained kinetic fractions before choosing any accuracy threshold.
Then choose a clipping rate $b_1$ and much larger thresholds $A_2$ and
$B\ge A_2$. For each path label call an active gap cheap when
\[
                  e_i=E(y_i)+\sigma_i\le A_2.
\]
There are boundedly many cheap gaps, and their total polynomial degree
is bounded in terms of the thresholds: a polynomial power also costs
at least one positive frequency. Their rates range in a fixed finite
set, independent of $K_0$ and $P$.

We first collect the absolute estimates with the cheap lengths set
to $s_0$. In this substitution the absolute polynomial coefficients
are evaluated at $s_0$; the binomial expansion is bounded by the same
polynomial-in-path-distance majorant used above. Removing or restoring
the bounded cheap shifts and their kinetic factors at $s_0$ costs only
a threshold-dependent constant. The unused half of the path decay
controls high shifts. A further fraction of the kinetic energy
controls high original energies. To apply the root stability estimate
in the presence of probes, compare the actual kinetic momentum with
the root momentum; their difference is bounded. Borrowing a slightly
larger fraction of actual energy loses at most a fixed constant per
unit length. The distant-path and bend losses are absorbed by the
bin-square penalty. On the two edges the remaining loss is therefore
\[
              e^{C(2\Delta_1+R_L+R_R)},
\]
with $C$ fixed independently of the thresholds and Taylor accuracy.

By taking $A_2,B$ large enough, every noncheap gap that is long or
has high energy retains a penalty
\[
                   e^{-(b_1+\lambda_2)d_i'},
\]
where $\lambda_2$ can be any prescribed large fixed number. The only
other noncheap gaps have $E(y_i)\le B$ and $d_i'<s_0$; there are
boundedly many and their total length is bounded. Reserve a small
additional kinetic fraction for Gaussian tails in $f$ when there is
no cheap gap. If a cheap gap is present, its bounded momentum implies
$|f|\le C_B+CP$.

Here is the root-count summation with its constants made explicit.
Sorting cancels identical-root factorials. Species/order choices cost
$a_*^P$ for one fixed $a_*$, while marking the bounded number of low gaps
and summing allowed $f$ costs at most $(1+P)^{C_B}$. Use
\[
       (1+P)^{C_B}\le C'_B2^P;
\]
only the prefactor depends on $B$. If $l$ roots lie in one edge,
its ordered volume, in a unit bin of $R_j$, is bounded by
$(\Delta_1+\lceil R_j\rceil+1)^l/l!$. Consequently
\[
 \sum_{l\ge0}\frac{(2a_*)^l(\Delta_1+\lceil R_j\rceil+1)^l}{l!}
 =\exp\{2a_*(\Delta_1+\lceil R_j\rceil+1)\}.
\]
Arbitrary fixed exponential factors in $R_j$ are integrable by
localization. Each expensive integrated free gap contributes
$1/\lambda_2$ after the $b_1$ part of its penalty is reserved.
The number of free roots and free gaps differs by at most one on each
arc. Thus all but boundedly many free roots are paid for by these
integrated expensive gaps. Choose $\lambda_2$
to absorb the fixed exponential bases. These estimates give
$C_J e^{C\Delta_1}$ for the remaining data, with one exponent $C$
independent of $J$. No accuracy-dependent constant has been raised
to the unbounded number of edge roots.

\paragraph{Convolution, then extension.}
Suppose each arc has a cheap gap. For a fixed monomial in its cheap
lengths, convolve the factors $(d_i')^{k_i}e^{-e_i d_i'}$ exactly,
with $d_i'\ge s_0$ and sum
\[
            \lambda_1-R_j-\sum_{\rm other}d_i'.
\]
On its feasible interval the result is a finite exponential-polynomial
in this sum. This follows by translating the lower endpoints and
partial fractions of the Laplace transform. Its rates are among the
$e_i$, and the coefficient bounds depend only on the fixed finite rate
and degree set. Retain rates at most $b_1$.

The bounds just collected now justify extending their remaining
coefficient integrals to the product of their individual ranges: retain
$d_i'<s_0$ for inactive low-energy gaps and $d_i'\ge s_0$ for active
gaps, but remove their coupled sum constraint. A retained rate contributes
at most $e^{b_1R_j+b_1\sum_{\rm other}d_i'}$, times a bounded-degree
polynomial, to these coefficients. The reserved penalties on other
gaps and the arbitrary exponential moments of $R_j$ make this
integrable. Put $S_j=R_j+\sum_{\rm other}d_i'$. Beyond the old feasibility
boundary $S_j\ge\lambda_1-O_B(1)$. After paying $e^{b_1S_j}$,
reserve a further factor $e^{-tS_j}$, with $t$ as large as needed,
by the choice of $\lambda_2$ and the exponential moments of $R_j$.
The inactive low gaps have bounded total length. Hence the tail is
\[
 e^{-t(\lambda_1-O_B(1))}\,C_J e^{C\Delta_1},
\]
up to bounded-degree polynomial factors, and has any prescribed
exponential accuracy in $\lambda_1$. The discarded convolution rates
have the corresponding bound
\[
       e^{-b_1\lambda_1+b_1R_j+b_1\sum_{\rm other}d_i'}
\]
on the feasible interval, again with only bounded-degree polynomial
factors. If an arc has no cheap gap, its length constraint gives this
smallness directly; eliminate a last gap there rather than convolving.
The other arc costs at most the already bounded coefficient integrals
and polynomial factors. Extending the earlier restriction
$V_1\le\lambda_1/2$ is justified by its retained superexponential
penalty.

Thus the coefficient integrals are absolutely summable, and their
errors have arbitrarily high exponential accuracy in $\lambda_1$ at
cost $C_J e^{C\Delta_1}$. Since
$\lambda_1=6\log s_1$ and $\Delta_1=6\log K_0$, this is the claimed
$C_JK_0^{C_2}s_1^{-J}$ form, with fixed $C_2$. The large clipping rate
has multiplied $R_j$, whose exponential moments change only $C_J$;
it has not multiplied $\Delta_1$.

\par\medskip\noindent\textbf{Minimal rates.}
What remains is polynomial powers in \(\lambda_1\) times rates \(e_L+e_R\) chosen among the cheap rates on the respective arcs (and the earlier nonnegative Taylor orders). Coefficients retain the stated bounds and holomorphy, since other positions, signs of real orders, angular parameters and phase classes in these prescriptions are fixed independent of \(\lambda_1\) and the composition in its neighborhood. Uniformity along the bounded imaginary shifts allows Cauchy averaging. For \(p_*=2\), between selected gaps in increasing real order across the center one gains \(3/2\) from the marks and only nonnegative other changes. Hence \(E(y_L)+E(y_R)\ge1/4\), giving \(\gamma\ge6/4\). At equality there is no Taylor correction and the two shifts vanish; the chosen momenta are \(-3/4,3/4\). No changing roots are between these gaps and the marks (uncoupled roots in the free region all change the momentum positively in increasing real order), so all other free gaps have strictly larger kinetic energy on their respective arcs. Thus both convolution poles are simple, with no contributing polynomial powers there since a path crossing either would shift its rate. For \(p_*=0\), momenta are in \(f_0+\mathbb Z\), giving \(6(E(y_L)+E(y_R))\ge1/6\), equality only for \(y_L=y_R=f_0\). For that order there can be no shifts or Taylor corrections and each convolution pole is simple without path powers by strict unit-or-larger changes along the relevant free gaps. Changing a momentum costs at least \(6/9\) in exponent, and positive shifts or Taylor powers at least \(1/2\). This proves \eqref{L:eq9}.


\end{proof}

\subsection{Physical interpolation}\label{L:step16}

First in the unmarked case the comparison \(|Q|\ge c_1{\cal D}\) forces nonvanishing of the leading coefficient, locally even at nonrational limiting compositions. Indeed fix any possible residue class of the family/angle counts occurring along an infinite sequence. To approach an admissible limit on it one can take the counts within bounded error of the exact scaled composition, preserving residues (and orientation at a tie if needed). One can use either nearly equal splitting since the identity computes the symmetric \(Q^2\) regardless of the chosen split. If the coefficient there vanished, holomorphy gives error \(O(1/|Z|)\) in evaluating it along this approximation, contradicting the comparison since \(5/24=3/8-1/6<1/2\). Compactness thus gives
\[
 |Q(Z)|^2\asymp {\cal D}(Z)^2 |Z|^{5/24}.
\]
Also the mass \(C_a\) of loops visiting just a prescribed mark (no constraint at other sites) is bounded above by a constant at any such physical list. Use the one-port mark reduction to \(E_+(Z;t)\) at \(t=a\) (take the seam there). By its formula and the cap reduction for \(Q(Z,q^3a)\) the nonconstant term is a bounded prefactor times \(Q(Z,a)Q(Z\setminus a)/Q(Z)^2\). The just-proved estimate gives the bound since the products in \({\cal D}\) cancel up to fixed factors. In particular \(0\le C_{ab}\le C_a\le C\).

For bounded $K_0$, equation~\eqref{L:eq9} now gives~\eqref{L:eq8}
directly. For large integral $K_0$, put $m_0=K_0^{1/4}$. We calibrate
the marked expansion using physical values near this smaller scale.
Choose one sufficiently large accuracy $J>4C_2+1$ and keep its finite
rate and degree list fixed. At $m_0$ the remainder satisfies
\[
 K_0^{C_2}m_0^{-J}=o(K_0^{-5/24})
                   =o((K_0m_0)^{-1/6}).
\]
The strict power margin also absorbs any fixed logarithmic factor.

\paragraph{Interpolating the angle fraction.}
Fix a target phase class and common physical fraction $\theta$.
At each of finitely many sizes $m'$ near fixed multiples of $m_0$,
choose an integer $m'$ within bounded error preserving the size class.
Use equally spaced neighboring admissible later counts $j'$ of the
prescribed class in the first family, and exactly $K_0j'$ in the second.
Choose them within bounded distance of $m'\theta$ and within the physical
range; one-sided choices suffice at an endpoint. The bound $C_{ab}\le C$
and the unmarked $Q$ estimate imply
\[
 |G_Z^{(2)}|\le C(K_0m_0)^{-1/6}
\]
at these physical points. Formula~\eqref{L:eq9} gives the same bound for
its truncated expression there.

Interpolate that expression in the common angle parameter. A fixed
sufficiently large number of the neighboring grid points has bounded
interpolation weights. Holomorphy and the coefficient bound give a
derivative remainder at most
\[
 C_JK_0^{C_2}(1+\log m_0)^{C_J}m_0^{-A}
\]
for any prescribed fixed interpolation order $A$. Choose $A$ large
enough that this and the formula errors are negligible. We have therefore
bounded the truncated expression at the exact target fraction $\theta$,
without a divisibility requirement on $m'\theta$.

\paragraph{Interpolating the size.}
Write $m=m_0e^u$. The truncated expression is a linear combination of
the finitely many functions $e^{-\gamma u}u^i$ in its fixed rate and
degree list. Choose fixed evaluation values of $u$ for which their
matrix is invertible; linear independence guarantees such a choice.
The bounded rounding of each $m'$ perturbs these values by $O(m_0^{-1})$,
so the inverses stay bounded. Thus every coefficient in this shifted
basis, including its factor $m_0^{-\gamma}$, is
$O((K_0m_0)^{-1/6})$ at the exact target composition.

For $u\ge0$, the higher-rate polynomial terms are bounded by a constant
times $e^{-3u/2}$. At the rate $3/2$ there is no polynomial factor.
Consequently
\[
 |G_Z^{(2)}|\le C(K_0m_0)^{-1/6}(m/m_0)^{-3/2}
                       +O_J(K_0^{C_2}m^{-J}),\qquad m\ge m_0.
\]
Undo the normalization, using $|Z|\asymp K_0m$ and the unmarked
$Q$ estimate. The principal term gives
\[
 C_{ab}\le C(K_0m)^{1/6}(K_0m_0)^{-1/6}(m/m_0)^{-3/2}
       =C m^{-4/3}K_0^{1/3}.
\]
The remainder relative to this target is at most
$C_JK_0^{C_2+5/24-J/4}$ for $J\ge3/2$, so the same choice of $J$
absorbs it. All choices can be made on fixed local neighborhoods;
a finite cover of the composition interval makes the constant uniform.
This proves~\eqref{L:eq8}.

\subsection{The polygon second moment}

We now convert the two-port estimate into a bound on all ordered vertex
pairs of a polygon. Good displacements are counted directly; polygons
concentrated near a lattice side line require a separate width argument.
For plane honeycomb polygons \(P\), modulo primitive lattice translations,
the resulting estimate is
\begin{equation}
 \sum_{R\le{\rm diam}\,P<2R} X^{|P|}|P|^2\le C R^{2/3}(1+\log\log R).\label{L:eq10}
\end{equation}
\paragraph{Pairs away from lattice side lines.}
Call ordered vertex pairs good for displacements of length at least \(R^{1/3}\) and at Euclidean distance greater than a sufficiently large fixed constant from all lattice side lines through zero. To count their weighted mass fix the first polygon vertex (either type) and sum over displacements to the other. A good displacement is well inside some rotated \(e_0,e_1\) cone. At each visited vertex the polygon uses two of three midports and therefore at least one oriented-side choice in these two directions. Sum over both choices there. For any prescribed such pair, one forward zigzag can include those two sides: between the end of the first and start of the second is a nonnegative integral combination of \(e_0,e_1\). If it takes \(m\) steps, \(m\) is comparable to the distance. On the rest of the zigzag period after the second mark take \(K_0\) copies of the intermediate step counts, \(K_0\asymp R/m\) integral and large enough that the period vector has length \(>4R\). Reflection exchanging the side directions allows the later proportion \(\le1/2\). Projection now embeds all the required polygons injectively at the marks (lift at the first mark), preserving physical weights. Thus \eqref{L:eq8} applies since \(m\gtrsim R^{1/3}\). At dyadic displacement scale \(r\) this costs \(R^{1/3} r^{1/3}\) up to constants. Hence \eqref{L:eq10} holds with the good-pair count in place of the squared length.

\paragraph{Polygons concentrated near a side line.}
We can discard with superpolynomial length-weighted accuracy polygons of transverse width below \(R/\log^2 R\) relative to any lattice side direction. Project these contractibly to a cylinder along a different (nonparallel) lattice side with period \(\lceil C R/\log^2 R\rceil\), \(C\) large fixed; translates do not meet by the width bound, whereas the height there is of order at least \(R\). The exponential pressure height estimate applies, also with length powers, and a plane polygon modulo translations is recovered by lifting.

For length \(l\gtrsim R\), short displacements account for \(O(lR^{2/3})\) pairs. If there is not already a fixed positive fraction good, some constant-width band through one vertex parallel to a lattice side then contains a positive fraction of all vertices by pigeonhole, implying \(l=O(R)\). Group these exceptional polygons by that side direction and dyadic transverse width \(w\gtrsim R/\log^2R\) relative to it. At least \(c w\) vertices lie at distance \(\gtrsim w\) from the band, by the width and bounded step size. For each, the other-direction side-line tubes through it meet the band in bounded regions, and a sufficiently small constant times \(R\) as minimum displacement length still leaves a positive fraction of the band vertices giving good partners (the band has constant width). Thus at least a fraction \(c' w/R\) of all ordered pairs are good with displacement of order \(R\), transverse component at most \(Cw\). There are only \(O(wR)\) such displacement vectors and each costs \(O(R^{-4/3})\) by the argument above. The good-pair mass in this bin is therefore at most
\[
 C(wR)R^{-4/3}=CwR^{-1/3}.
\]
Each such polygon has at least $c(w/R)|P|^2$ good pairs, so its
squared-length mass, summed over the bin, is at most
\[
 C\frac{R}{w}\,wR^{-1/3}=CR^{2/3}.
\]
Only $O(1+\log\log R)$ width bins lie between $R/\log^2R$ and $O(R)$.
Together with the ordinary good-pair estimate and the negligible thin
polygons, this proves~\eqref{L:eq10}.

